\documentclass[12pt,a4paper]{article}

% 基础包
\usepackage[utf8]{inputenc}
\usepackage{xeCJK}
\usepackage{fontspec}
\usepackage{geometry}
\usepackage{graphicx}
\usepackage{xcolor}
\usepackage{amsmath}
\usepackage{amssymb}
\usepackage{listings}
\usepackage{tikz}
\usepackage{booktabs}
\usepackage{array}
\usepackage{enumitem}
\usepackage{fancyhdr}
\usepackage{titlesec}
\usepackage{tcolorbox}
\usepackage{tabularx}
\usepackage{ragged2e}
\usepackage{draftwatermark}
\usepackage{eso-pic}

\setCJKmainfont{SimSun}
\setCJKsansfont{SimHei}
\setCJKmonofont{FangSong}
\newfontfamily\monoascii{Courier New}

\geometry{margin=2.5cm,top=3cm,bottom=2.5cm}

% 颜色定义
\definecolor{PrimaryBlue}{HTML}{2E5090}
\definecolor{AccentOrange}{HTML}{E67E22}
\definecolor{DarkText}{HTML}{2C3E50}
\definecolor{LightBG}{HTML}{F8F9FA}
\definecolor{ExamBG}{HTML}{FFF8E7}
\definecolor{TheoryBG}{HTML}{F0F7FF}
\definecolor{ConfidentialRed}{HTML}{C41E3A}
\definecolor{InternalGray}{HTML}{4A4A4A}

% tcolorbox设置
\tcbuselibrary{skins,breakable}

\newtcolorbox{keypointbox}{
  colback=LightBG,
  colframe=PrimaryBlue,
  boxrule=0.5pt,
  arc=4pt,
  fonttitle=\bfseries\color{PrimaryBlue},
  title=重点提示,
  left=8pt,right=8pt,top=8pt,bottom=8pt
}

\newtcolorbox{exambox}{
  colback=ExamBG,
  colframe=AccentOrange,
  boxrule=0.5pt,
  arc=4pt,
  fonttitle=\bfseries\color{AccentOrange},
  title=重点提示,
  left=8pt,right=8pt,top=8pt,bottom=8pt
}

\newtcolorbox{theorybox}{
  colback=TheoryBG,
  colframe=PrimaryBlue,
  boxrule=0.5pt,
  arc=4pt,
  fonttitle=\bfseries\color{PrimaryBlue},
  title=理论解析,
  left=8pt,right=8pt,top=8pt,bottom=8pt,
  before upper={\setlength{\parindent}{0pt}},
  after upper={\setlength{\parindent}{0pt}\normalfont}
}

% 代码样式
\lstset{
  basicstyle=\small\monoascii,
  breaklines=true,
  frame=single,
  backgroundcolor=\color{LightBG},
  rulecolor=\color{PrimaryBlue!30},
  numbers=left,
  numberstyle=\tiny\color{gray},
  keywordstyle=\color{PrimaryBlue}\bfseries,
  commentstyle=\color{gray}\itshape,
  stringstyle=\color{AccentOrange},
  showstringspaces=false,
  tabsize=2,
  xleftmargin=2em,
  framexleftmargin=1.5em
}

% Header/footer - 内部资料样式
\fancyhf{}
\fancyhead[L]{\small\color{ConfidentialRed}\bfseries [内部资料·禁止外泄]}
\fancyhead[R]{\small\color{DarkText}2026年安徽省区块链技术及应用技能赛}
\fancyfoot[C]{\thepage}
\fancyfoot[L]{\scriptsize\color{ConfidentialRed} 内部培训资料，严禁外传}
\fancyfoot[R]{\scriptsize\color{InternalGray} 仅供参赛选手使用}
\addtolength{\headheight}{4pt}
\pagestyle{fancy}

% 添加水印
\SetWatermarkText{内部资料}
\SetWatermarkScale{0.8}
\SetWatermarkColor{ConfidentialRed}
\SetWatermarkAngle{45}
\SetWatermarkLightness{0.9}

% 标题格式
\titleformat{\section}{\Large\bfseries\color{PrimaryBlue}}{\thesection}{1em}{}
\titleformat{\subsection}{\large\bfseries\color{DarkText}}{\thesubsection}{1em}{}
\titleformat{\subsubsection}{\normalsize\bfseries\color{DarkText}}{\thesubsubsection}{1em}{}

\begin{document}

% =================== 封面 ===================
\begin{titlepage}
\centering

\vspace*{2cm}

{\Huge\bfseries\color{PrimaryBlue} 区块链技术及应用技能赛\par}
\vspace{0.8cm}
{\LARGE\color{DarkText} 赛前培训资料\par}

\vspace{2cm}

\begin{tikzpicture}[scale=1.2]
\draw[PrimaryBlue, line width=1pt] (-4,0) -- (4,0);
\draw[PrimaryBlue, fill=PrimaryBlue] (-4,0) circle (4pt);
\draw[PrimaryBlue, fill=PrimaryBlue] (-2,0) circle (4pt);
\draw[PrimaryBlue, fill=PrimaryBlue] (0,0) circle (4pt);
\draw[PrimaryBlue, fill=PrimaryBlue] (2,0) circle (4pt);
\draw[PrimaryBlue, fill=PrimaryBlue] (4,0) circle (4pt);
\draw[PrimaryBlue, thick] (-4,0) -- (-2,0);
\draw[PrimaryBlue, thick] (-2,0) -- (0,0);
\draw[PrimaryBlue, thick] (0,0) -- (2,0);
\draw[PrimaryBlue, thick] (2,0) -- (4,0);

\node at (-4,-0.6) {\small 区块1};
\node at (-2,-0.6) {\small 区块2};
\node at (0,-0.6) {\small 区块3};
\node at (2,-0.6) {\small 区块4};
\node at (4,-0.6) {\small 区块5};
\end{tikzpicture}

\vspace{1.5cm}

{\Large\color{DarkText} 参赛选手内部培训资料\par}
\vspace{0.5cm}
{\large\color{InternalGray} 2026年安徽省区块链技术及应用技能赛\par}

\vspace{1.2cm}

% 内部资料标识框
\begin{tcolorbox}[
  colback=ConfidentialRed!5,
  colframe=ConfidentialRed,
  boxrule=1.5pt,
  arc=3pt,
  width=0.75\textwidth,
  halign=center
]
{\Huge\bfseries\color{ConfidentialRed} 内部资料}\\[8pt]
{\Large\color{InternalGray} 严禁外传 · 禁止复制 · 仅供参赛选手使用}
\end{tcolorbox}

\vspace{2.5cm}

{\large\color{InternalGray} 
\begin{tabular}{c}
安徽省电子学会 · 安徽省计算机学会 · 巢湖学院 · 中科晶格 \\
\vspace{0.3cm}
2026年5月
\end{tabular}
}

\vfill

% 保密声明框
\begin{tcolorbox}[
  colback=ConfidentialRed!8,
  colframe=ConfidentialRed,
  boxrule=2pt,
  arc=4pt,
  width=0.85\textwidth,
  halign=center
]
{\Large\bfseries\color{ConfidentialRed} 保密声明}\\[8pt]
{\normalsize\color{InternalGray}
本资料为\textbf{内部培训资料}，仅供参赛选手使用。\\[3pt]
\textbf{严禁}以任何形式外传、复制、拍照或上传至互联网。\\[3pt]
违者将取消参赛资格，并追究相关责任。
}
\end{tcolorbox}

\end{titlepage}

\newpage
\tableofcontents

\newpage

% =================== 第1章 区块链零基础入门 ===================
\section{区块链零基础入门}

\subsection{计算机基础知识}

本节梳理最核心的计算机基础概念，可能以选择题或判断题形式出现。

\subsubsection{进制与编码}

\begin{itemize}[leftmargin=1.5em]
\item \textbf{二进制（Bin）}：计算机内部只用0和1两种状态，比特币的Hash运算是二进制层面的计算
\item \textbf{十六进制（Hex）}：二进制太长，用十六进制表示更简洁。SHA-256的输出就是64个十六进制字符
\item \textbf{字符编码}：UTF-8是互联网上最通用的字符编码，中文在区块链系统中以UTF-8存储
\item \textbf{Base58编码}：比特币地址使用Base58Check编码，去掉了容易混淆的字符（0、O、l、I），更 human-readable
\end{itemize}

\begin{keypointbox}
记住：SHA-256输出64个十六进制字符 = 256位二进制 = 32字节
\end{keypointbox}

\subsubsection{网络基础}

\begin{itemize}[leftmargin=1.5em]
\item \textbf{IP地址}：互联网上的"门牌号"，区块链节点通过IP互相通信
\item \textbf{DNS}：域名系统，把人类可读的网址（如google.com）转成IP地址
\item \textbf{HTTP/HTTPS}：网页传输协议，区块链浏览器使用HTTPS保证安全
\item \textbf{P2P网络}：点对点网络，区块链的基础架构。没有中心服务器，所有节点平等
\end{itemize}

\subsubsection{数据库基础}

\begin{itemize}[leftmargin=1.5em]
\item \textbf{关系型数据库}：MySQL、Oracle，用表格存储数据，有行和列
\item \textbf{键值对数据库}：Redis、LevelDB，用Key-Value存储，区块链常用
\item \textbf{分布式数据库}：数据分散存储在多台机器上，区块链是特殊的分布式数据库
\item \textbf{ACID特性}：原子性、一致性、隔离性、持久性。传统数据库追求ACID，区块链追求最终一致性
\end{itemize}

\subsubsection{JSON格式}

区块链中大量数据以JSON格式传输和存储。

\begin{lstlisting}
{
  "name": "Alice",
  "age": 25,
  "isStudent": false,
  "courses": ["Math", "CS", "Blockchain"]
}
\end{lstlisting}

重点提示：JSON格式看起来像这样——``{\"key\": \"value\", \"number\": 123}''。以太坊的智能合约调用和大多数区块链浏览器都返回JSON数据。

\subsection{职业道德与法律法规}

\subsubsection{职业道德基本知识}

\begin{itemize}[leftmargin=1.5em]
\item \textbf{诚实守信}：区块链行业强调透明和可追溯，从业者必须诚实记录交易和数据
\item \textbf{保护用户隐私}：虽然区块链数据公开，但用户的真实身份需要保护，不能随意泄露
\item \textbf{持续学习}：区块链技术更新快，从业者需要不断学习新技术和新标准
\item \textbf{安全意识}：私钥管理、智能合约安全是区块链的核心，任何疏忽都可能导致巨大损失
\end{itemize}

\subsubsection{相关法律法规}

\begin{itemize}[leftmargin=1.5em]
\item \textbf{《中华人民共和国密码法》}（2020年1月1日实施）
  \begin{itemize}
  \item 规范密码应用和管理，促进密码事业发展
  \item 密码分为核心密码、普通密码和商用密码
  \item 区块链中使用的SM2/SM3/SM4属于商用密码
  \end{itemize}
\item \textbf{《区块链信息服务管理规定》}（2019年2月15日实施）
  \begin{itemize}
  \item 要求区块链信息服务提供者进行备案
  \item 禁止利用区块链传播违法信息
  \item 要求对使用者进行真实身份信息认证
  \end{itemize}
\item \textbf{《数据安全法》}（2021年9月1日实施）
  \begin{itemize}
  \item 规范数据处理活动，保障数据安全
  \item 重要数据需要分类分级保护
  \end{itemize}
\item \textbf{《个人信息保护法》}（2021年11月1日实施）
  \begin{itemize}
  \item 规范个人信息处理活动
  \item 区块链上存储个人信息需要符合法律规定
  \end{itemize}
\end{itemize}

\begin{exambox}
重点提示：法律部分需要知道有哪些规定存在，以及规定的核心目的（如密码法是为了规范商业密码使用）。不需要背诵条文。
\end{exambox}

\subsection{通证与资产应用（补充）}

通证（Token）是区块链的独特概念。

\subsubsection{什么是通证（Token）？}

简单说：\textbf{通证就是记录在区块链上的"权益凭证"}。

\begin{theorybox}
\textbf{通证的三个特征：}
\begin{enumerate}
\item \textbf{可流通}：可以在区块链上自由转让（如转让NFT所有权）
\item \textbf{可验证}：任何人都可以在链上验证其真伪（查区块链即可）
\item \textbf{可编程}：可以通过智能合约设定规则（自动分红、限量发行等）
\end{enumerate}
\end{theorybox}

类比：门票就是一种通证——有编号（可验证）、可转让（给他人）、有时效（可编程）。

\subsubsection{通证的分类}

\begin{itemize}[leftmargin=1.5em]
\item \textbf{同质化通证（FT，Fungible Token）}：每个Token都一样，可以互换。类比人民币——每张100元价值相等。代表：比特币（BTC）、以太坊（ETH）
\item \textbf{非同质化通证（NFT，Non-Fungible Token）}：每个Token唯一，不可互换。类比房产证——每个房产证对应一套唯一的房子。代表：数字艺术品、游戏道具
\item \textbf{证券化通证（STO）}：代表真实资产的所有权（如房地产、股票），受证券法规管
\end{itemize}

\begin{keypointbox}
重点：BTC/ETH = 同质化通证（FT），每个都一样；NFT = 非同质化通证，每个都唯一。
\end{keypointbox}

% =================== 第2章 区块链核心概念 ===================
\section{区块链核心概念}

\subsection{什么是区块链？}

区块链（Blockchain）本质上是一个\textbf{分布式数据库}，但它和普通数据库有一个本质区别：\textbf{数据一旦写入，就几乎不可能被修改或删除}。

用更通俗的话说：\textbf{区块链就是一个大家共同维护的、无法作弊的电子账本}。

\subsubsection*{从记账场景理解}

传统模式：你把钱存银行，银行在它的服务器上记账。你只能\textbf{相信}银行不会篡改你的记录。

区块链模式：\textbf{所有参与者都保存一份完整的账本副本}。任何人修改了自己的账本，其他人的账本都会拒绝匹配。就像全班同学都抄了一份课程表，有人偷偷改了自己的，老师一比对就知道谁在说谎。

\begin{theorybox}
\textbf{区块链的核心特点：}
\begin{enumerate}
\item \textbf{去中心化}：没有单一的控制者，所有节点平等
\item \textbf{不可篡改}：数据一旦写入，修改成本极高（需要控制超过51\%的节点）
\item \textbf{可追溯}：每一笔交易都有完整的历史记录，可以追踪资金流向
\item \textbf{透明性}：所有交易对参与者公开（虽然身份可能是匿名的）
\end{enumerate}
\end{theorybox}

\subsection{区块链 vs 传统数据库}

\begin{center}
\begin{tabular}{|l|p{5cm}|p{5cm}|}
\hline
\rowcolor{LightBG}
\textbf{特性} & \textbf{传统数据库} & \textbf{区块链} \\
\hline
数据修改 & 管理员可以随时修改 & 几乎不可能修改 \\
\hline
控制权 & 中心化（公司/机构控制） & 去中心化（共识机制决定） \\
\hline
信任基础 & 相信中心机构 & 相信密码学和共识机制 \\
\hline
查询速度 & 快（毫秒级） & 慢（秒级甚至分钟级） \\
\hline
存储成本 & 低 & 高（每个节点都存一份） \\
\hline
适合场景 & 高频交易、大数据 & 价值转移、存证溯源 \\
\hline
\end{tabular}
\end{center}

\begin{keypointbox}
区块链不是万能的！它牺牲了效率换取了信任。如果你的场景不需要"去信任"，传统数据库可能更合适。
\end{keypointbox}

\subsection{区块链的类型}

\begin{itemize}[leftmargin=1.5em]
\item \textbf{公链（Public Blockchain）}：任何人都可以参与，完全去中心化
  \begin{itemize}
  \item 代表：比特币、以太坊
  \item 特点：匿名、开放、抗审查
  \item 共识：PoW、PoS
  \end{itemize}
\item \textbf{联盟链（Consortium Blockchain）}：多个机构共同维护，部分去中心化
  \begin{itemize}
  \item 代表：Hyperledger Fabric、R3 Corda、FISCO BCOS
  \item 特点：需要许可、效率较高、适合企业
  \item 共识：PBFT、Raft
  \end{itemize}
\item \textbf{私链（Private Blockchain）}：单一机构控制，中心化
  \begin{itemize}
  \item 代表：企业内部系统
  \item 特点：完全可控、效率最高、去中心化程度低
  \item 用途：内部审计、数据共享
  \end{itemize}
\end{itemize}

\begin{exambox}
中国主要发展\textbf{联盟链}，因为：1）符合监管要求；2）性能更高；3）适合企业级应用。比特币、以太坊是公链的代表。
\end{exambox}

% =================== 第3章 密码学基础 ===================
\section{密码学基础}

密码学是区块链的基石。区块链中几乎所有的安全机制都依赖于密码学。本章内容占比很高，务必理解原理。

\subsection{Hash函数（哈希函数）}

\subsubsection{什么是Hash？}

Hash函数是一个\textbf{单向函数}，可以把任意长度的输入转换成固定长度的输出。

\begin{lstlisting}
输入: "Hello World"
SHA-256输出: a591a6d40bf420404a011733cfb7b190d62c65bf0bcda32b57b277d9ad9f146e

输入: "Hello World!" (只多了一个感叹号)
SHA-256输出: 7f83b1657ff1fc53b92dc18148a1d65dfc2d4b1fa3d677284addd200126d9069
\end{lstlisting}

\begin{theorybox}
\textbf{Hash函数的核心特性：}
\begin{enumerate}
\item \textbf{单向性}：知道输出，无法反推出输入（除非暴力破解）
\item \textbf{确定性}：同样的输入，永远得到同样的输出
\item \textbf{雪崩效应}：输入哪怕改变1bit，输出完全不同
\item \textbf{抗碰撞}：很难找到两个不同的输入，产生相同的输出
\end{enumerate}
\end{theorybox}

\subsubsection{区块链中的Hash应用}

\begin{enumerate}[leftmargin=1.5em]
\item \textbf{区块链接}：每个区块存储前一个区块的Hash，形成链式结构
\item \textbf{交易指纹}：每笔交易计算Hash作为唯一标识
\item \textbf{数据完整性}：任何数据篡改都会导致Hash变化，立即被发现
\item \textbf{工作量证明（PoW）}：矿工通过计算Hash来竞争记账权
\end{enumerate}

\subsubsection{常用Hash算法}

\begin{itemize}[leftmargin=1.5em]
\item \textbf{SHA-256}：比特币使用的算法，输出256位（64个十六进制字符）
\item \textbf{SM3}：中国国密标准，输出256位，与SHA-256安全强度相当
\item \textbf{Keccak-256}：以太坊使用的算法（SHA-3标准）
\item \textbf{RIPEMD-160}：比特币地址生成中使用的算法，输出160位
\end{itemize}

\begin{keypointbox}
重点：SHA-256输出是64个十六进制字符（每字符=4bit，$64\times4=256$bit）。SM3是中国国密标准，与SHA-256安全强度相当。区块链中优先学习SHA-256和SM3。
\end{keypointbox}

\subsubsection{Python计算Hash}

\begin{lstlisting}
import hashlib

# SHA-256计算
data = b"Hello Blockchain"
hash_result = hashlib.sha256(data).hexdigest()
print(f"SHA-256: {hash_result}")

# 验证雪崩效应
data1 = b"blockchain"
data2 = b"blockchaiN"  # 只改一个大写字母
h1 = hashlib.sha256(data1).hexdigest()
h2 = hashlib.sha256(data2).hexdigest()
print(f"原始: {h1}")
print(f"修改: {h2}")
print(f"不同字符数: {sum(c1 != c2 for c1, c2 in zip(h1, h2))}")
\end{lstlisting}

\subsection{非对称加密（公钥加密）}

\subsubsection{核心概念}

非对称加密使用\textbf{一对密钥}：
\begin{itemize}[leftmargin=1.5em]
\item \textbf{公钥（Public Key）}：可以公开分享，用于加密和验证签名
\item \textbf{私钥（Private Key）}：必须保密，用于解密和创建签名
\end{itemize}

\begin{theorybox}
\textbf{非对称加密的两个用途：}
\begin{enumerate}
\item \textbf{加密通信}：A用B的公钥加密，只有B的私钥能解密
\item \textbf{数字签名}：A用私钥签名，任何人可以用A的公钥验证
\end{enumerate}
\end{theorybox}

\subsubsection{椭圆曲线加密（ECC）}

区块链主要使用\textbf{椭圆曲线加密}（ECC），而不是RSA。

\begin{itemize}[leftmargin=1.5em]
\item \textbf{优势}：相同安全强度下，ECC的密钥更短、计算更快
  \begin{itemize}
  \item 256位ECC $\approx$ 3072位RSA的安全强度
  \item 比特币、以太坊使用secp256k1曲线
  \item 中国国密SM2使用sm2p256v1曲线
  \end{itemize}
\item \textbf{比特币的密钥生成}：
  \begin{enumerate}
  \item 生成256位随机数作为私钥
  \item 使用ECC点乘法从私钥生成公钥
  \item 对公钥进行Hash运算生成地址
  \end{enumerate}
\end{itemize}

\begin{keypointbox}
\textbf{重点：}256位ECC $\approx$ 3072位RSA的安全强度。区块链（比特币、以太坊）使用ECC（secp256k1曲线）。中国使用SM2（sm2p256v1曲线）。公钥可以分享，私钥必须保密。
\end{keypointbox}

\subsection{数字签名}

数字签名是区块链身份认证的核心机制。

\subsubsection{签名过程}

\begin{enumerate}[leftmargin=1.5em]
\item \textbf{创建签名}：用私钥对消息Hash进行"加密"
\item \textbf{验证签名}：用公钥"解密"并比对Hash
\end{enumerate}

\begin{lstlisting}
签名流程：
消息 -> Hash函数 -> Hash值 -> 私钥签名 -> 签名结果
验证流程：
签名结果 + 公钥 -> 解密出Hash值 -> 与消息的新Hash比对
\end{lstlisting}

\begin{keypointbox}
\textbf{重点：}
\begin{itemize}
\item 签名 = 私钥加密（Hash值）
\item 验证 = 公钥解密（比对Hash）
\item 只有私钥持有者能创建有效签名
\item 任何人都能验证签名真伪
\item 签名不可伪造、不可抵赖
\end{itemize}
\end{keypointbox}

\subsubsection{区块链中的签名应用}

\begin{itemize}[leftmargin=1.5em]
\item \textbf{交易签名}：证明交易确实由资产所有者发起
\item \textbf{区块签名}：证明区块确实由矿工/验证者生成
\item \textbf{智能合约调用}：证明调用者身份
\end{itemize}

\subsection{Merkle Tree（默克尔树）}

Merkle Tree是区块链中组织交易数据的核心数据结构。

\subsubsection{结构和原理}

\begin{lstlisting}
交易列表: [Tx1, Tx2, Tx3, Tx4]

        [Root Hash]           <- Merkle Root（存区块头）
         /       \
   [Hash01]     [Hash23]
    /    \\       /    \
[Tx1]  [Tx2]  [Tx3]  [Tx4]   <- 叶子节点（交易Hash）

计算过程:
Hash0 = Hash(Tx1)
Hash1 = Hash(Tx2)
Hash2 = Hash(Tx3)
Hash3 = Hash(Tx4)
Hash01 = Hash(Hash0 + Hash1)
Hash23 = Hash(Hash2 + Hash3)
Root = Hash(Hash01 + Hash23)
\end{lstlisting}

\subsubsection{Merkle Tree的优势}

\begin{enumerate}[leftmargin=1.5em]
\item \textbf{快速验证}：验证一笔交易是否存在，只需要 $O(\log n)$ 的Hash计算
\item \textbf{轻节点友好}：SPV（简单支付验证）节点只需下载区块头（80字节），不需要下载所有交易
\item \textbf{防篡改}：任何一笔交易被修改，Merkle Root就会改变，立即被发现
\end{enumerate}

\begin{lstlisting}
验证 Tx2 是否在区块中：
路径：Tx2 -> Hash(Tx3) -> Hash0 -> ... -> Merkle Root
如果最终计算得到的Merkle Root与区块头中的一致 -> 交易存在
\end{lstlisting}

\begin{keypointbox}
重点：Merkle Root存储在区块头中。如果交易列表中有任何一笔被修改，其Merkle Root就会改变。注意Merkle Tree是二叉树，节点数是偶数（奇数时复制最后一个节点）。
\end{keypointbox}

% =================== 第4章 区块链底层数据结构 ===================
\section{区块链底层数据结构}

\subsection{区块结构}

一个区块由两部分组成：\textbf{区块头（Block Header）}和\textbf{区块体（Block Body）}。

\subsubsection{区块头（Block Header）}

固定80字节，包含6个字段：

\begin{center}
\begin{tabular}{|l|l|l|}
\hline
\rowcolor{LightBG}
\textbf{字段} & \textbf{长度} & \textbf{说明} \\
\hline
Version & 4字节 & 区块版本号 \\
\hline
Previous Block Hash & 32字节 & 前一个区块的Hash（链式结构） \\
\hline
Merkle Root & 32字节 & 所有交易的Merkle Tree根Hash \\
\hline
Timestamp & 4字节 & 区块创建时间（Unix时间戳） \\
\hline
Bits & 4字节 & 难度目标（编码后的目标值） \\
\hline
Nonce & 4字节 & 随机数（PoW挖矿用） \\
\hline
\end{tabular}
\end{center}

\begin{keypointbox}
重点：区块头总共80字节，是PoW工作量证明的直接操作对象。父区块哈希将区块串成链。Nonce是4字节无符号整数，意味着在4字节范围内（大约43亿次）找不到有效Hash时，需要其他策略（如改变Coinbase交易中的ExtraNonce）。
\end{keypointbox}

\subsubsection{区块体（Block Body）}

包含交易列表，大小不固定：
\begin{itemize}[leftmargin=1.5em]
\item 交易数量（变长整数）
\item 交易列表（每笔交易200-500字节不等）
\end{itemize}

比特币区块大小限制：1MB（隔离见证后4MB）

\subsection{交易结构（UTXO模型）}

比特币使用\textbf{UTXO（Unspent Transaction Output）}模型，不是账户余额模型。

\subsubsection{什么是UTXO？}

UTXO = \textbf{未花费的交易输出}。可以理解为"区块链上的零钱"。

\begin{lstlisting}
类比现实：比特币的"余额"就像你钱包里的现金。
  - 你有：1张100元、2张50元、5张10元
  - 钱包余额 = 100 + 50×2 + 10×5 = 250元

比特币也一样：
  - UTXO #A: 0.5 BTC（来自某笔交易）
  - UTXO #B: 1.2 BTC（来自另一笔交易）
钱包余额 = 1.7 BTC
\end{lstlisting}

\subsubsection{交易过程}

\begin{lstlisting}
小明要转1 BTC给小红：

小明的UTXO：
  - UTXO #1: 0.6 BTC
  - UTXO #2: 0.8 BTC
  - 总计：1.4 BTC

交易创建：
  输入：引用 UTXO #1 和 UTXO #2（作为花费）
  输出：
  - 新建 UTXO #A: 1.0 BTC（归属小红）
  - 新建 UTXO #B: 0.399 BTC（归属小明，找零）
  - 手续费：0.001 BTC（给矿工）

结果：
  - UTXO #1 和 #2 被标记为"已花费"
  - UTXO #A 和 #B 成为新的"未花费输出"
\end{lstlisting}

\begin{keypointbox}
重点：UTXO模型的核心是"输入引用过去的输出，产生新的输出"。理解UTXO vs 账户余额模型的区别：银行是"记账"（A账户减，B账户加），比特币是"记录每笔零钱的归属"（UTXO的流转）。以太坊使用的是账户余额模型。
\end{keypointbox}

% =================== 第5章 共识算法 ===================
\section{共识算法}

共识算法解决的是\textbf{分布式系统中如何达成一致}的问题。

\subsection{拜占庭将军问题}

由Leslie Lamport等人于1982年提出，是区块链共识机制最重要的理论来源。

\begin{lstlisting}
情境：10位拜占庭将军分别带领一支军队包围了一座敌人城市。

问题：
- 将军们只能通过信使传递消息（没有电话、没有网络）
- 其中一些将军可能是叛徒，会发送错误信息
- 忠诚的将军们需要就"进攻还是撤退"达成一致

问题：如何让所有忠诚的将军就"进攻还是撤退"达成一致的行动计划？
\end{lstlisting}

将将军视为区块链节点，信使视为网络消息，叛徒视为恶意节点。

\begin{theorybox}
\textbf{拜占庭容错（BFT）}：系统能容忍最多$f$个恶意节点，仍能达成一致。

\begin{itemize}
\item $n \geq 3f + 1$：系统能容忍最多$f$个恶意节点（即恶意节点不超过总数的1/3）
\item 10个将军，最多容忍3个叛徒
\item 100个节点，最多容忍33个恶意节点
\end{itemize}
\end{theorybox}

\subsection{PoW — 工作量证明（Proof of Work）}

\subsubsection{核心思想}

\textbf{谁算力强，谁有话语权}。通过消耗计算资源（算力）来竞争记账权。

矿工不断改变一个叫Nonce的随机数，对区块头反复计算Hash，直到Hash值小于某个目标值。

\begin{lstlisting}
目标：Hash(区块头 + Nonce) $<$\texttt{0000FFFFFFFFFFFFFFFF}\\ \texttt{FFFFFFFFFFFFFFFFFFFFFFFFFFFF}

矿工尝试：
Nonce=0    -> Hash=8f8e2a1b3c4d5e6f...
Nonce=1    -> Hash=3a4b5c6d7e8f9a0b...
Nonce=...（数十亿次尝试）
Nonce=12345678 -> Hash=00001a2b3c4d5e6f...  ✓ 找到啦！
\end{lstlisting}

\subsubsection{PoW的特点}

\begin{itemize}[leftmargin=1.5em]
\item \textbf{优点}：
  \begin{itemize}
  \item 安全性高：攻击成本极高（需要控制51\%算力）
  \item 去中心化程度高：任何人都可以参与挖矿
  \item 公平性：算力即权力，规则透明
  \end{itemize}
\item \textbf{缺点}：
  \begin{itemize}
  \item 能耗巨大：比特币年耗电量超过阿根廷全国
  \item 效率低：每秒只能处理3-7笔交易
  \item 算力集中：矿池控制大部分算力，有中心化趋势
  \end{itemize}
\end{itemize}

\begin{keypointbox}
重点：PoW的核心是"算力竞争，谁先找到有效Hash谁赢"。难度由Bits字段编码。51\%攻击是PoW的最大威胁，但实际攻击成本远超收益。比特币使用SHA-256d（两次SHA-256）作为PoW函数。
\end{keypointbox}

\subsection{PoS — 权益证明（Proof of Stake）}

\subsubsection{核心思想}

根据持有的代币数量和时间来选择区块验证者。持有的越多、被锁定的时间越长，被选中出块的概率越高。类似于"股东大会投票"，股份越多，话语权越大。

\begin{lstlisting}
PoS选择公式（简化）：
选择概率 $\propto$ 质押代币数量 × 质押时间

示例：
  - 持有1000 ETH，锁定1年 -> 权重 = 1000 × 1 = 1000
  - 持有500 ETH，锁定2年 -> 权重 = 500 × 2 = 1000
  - 权重相同，随机选一个
\end{lstlisting}

以太坊2.0的PoS使用\textbf{随机化（Randomized）}方法——RANDAO + VDF（可验证延迟函数）产生随机数，结合质押权重选择出块者。

\subsubsection{PoS的特点}

\begin{itemize}[leftmargin=1.5em]
\item \textbf{优点}：
  \begin{itemize}
  \item 能耗极低：比PoW节能99.95\%
  \item 效率高：出块速度快，每秒可处理数千笔交易
  \item 攻击成本高：攻击者需要购买并质押大量代币，攻击后代币会贬值
  \end{itemize}
\item \textbf{缺点}：
  \begin{itemize}
  \item 富者愈富：持币多的人更容易获得奖励
  \item 无利害关系（Nothing at Stake）：验证者可能在多个分叉上同时签名
  \item 初始分配问题：代币初始分配可能不公平
  \end{itemize}
\end{itemize}

\begin{keypointbox}
重点：PoS中"质押"（Stake）意味着将代币作为保证金。如果验证者作恶，其质押的代币会被罚没（Slashing）。以太坊PoS最多可罚没验证者32 ETH。这比PoW的51\%攻击更经济高效，能耗降低约99.95\%。
\end{keypointbox}

\subsection{DPoS — 委托权益证明}

\subsubsection{核心思想}

代币持有者\textbf{投票选举}出少数节点（如21个）负责出块和验证。类似于"代议制民主"，大家投票选代表。

\begin{lstlisting}
EOS的DPoS机制：
1. EOS持有者投票选出21个"区块生产者"（BP）
2. 21个BP轮流出块（每3秒一个区块）
3. BP获得区块奖励，需要回馈社区
4. 如果BP表现不好，持有者可以随时投票罢免
\end{lstlisting}

\subsubsection{DPoS的特点}

\begin{itemize}[leftmargin=1.5em]
\item \textbf{优点}：
  \begin{itemize}
  \item 效率极高：每秒可处理数千甚至上万笔交易
  \item 能耗极低：不需要大规模算力竞争
  \end{itemize}
\item \textbf{缺点}：
  \begin{itemize}
  \item 中心化程度高：只有21个节点控制网络
  \item 投票率低：大多数代币持有者不参与投票
  \item 去中心化程度降低——节点数量少，更容易被操控
  \end{itemize}
\end{itemize}

\subsection{PBFT — 实用拜占庭容错}

\subsubsection{核心思想}

通过\textbf{多轮投票}达成共识，只要恶意节点不超过总数的1/3，系统就能正常工作。

\subsubsection{三阶段协议}

\begin{enumerate}[leftmargin=1.5em]
\item \textbf{Request}：客户端向主节点发送请求
\item \textbf{Pre-Prepare}：主节点验证请求，广播给所有节点
\item \textbf{Prepare}：节点验证并广播Prepare消息
\item \textbf{Commit}：节点收到足够多Prepare后，广播Commit消息
\item \textbf{Reply}：节点收到足够多Commit后，执行并回复客户端
\end{enumerate}

\begin{theorybox}
\textbf{PBFT的容错能力：}
\begin{itemize}
\item 总节点数 $n$，恶意节点数 $f$
\item 系统能正常工作的条件：$n \geq 3f + 1$
\item 即：恶意节点不超过总数的1/3
\item 4个节点可以容忍1个恶意节点
\item 10个节点可以容忍3个恶意节点
\end{itemize}
\end{theorybox}

\subsubsection{PBFT的特点}

\begin{itemize}[leftmargin=1.5em]
\item \textbf{优点}：
  \begin{itemize}
  \item 即时确认：交易一旦写入就最终确定，不需要等待多个区块确认
  \item 高吞吐量：每秒可处理数千笔交易
  \item 确定性：不会出现分叉
  \end{itemize}
\item \textbf{缺点}：
  \begin{itemize}
  \item 扩展性差：节点数量增加，通信复杂度$O(n^2)$急剧上升
  \item 中心化趋势：通常用于联盟链，节点数量有限
  \item 不适合公链：公链节点太多，PBFT通信开销太大
  \end{itemize}
\end{itemize}

\begin{keypointbox}
重点：
\begin{itemize}
\item PBFT适合联盟链（节点少、信任度较高）
\item PoW/PoS适合公链（节点多、无需许可）
\item PBFT是"投票"机制，PoW是"算力竞争"机制
\item 联盟链常用PBFT（如Hyperledger Fabric），公链用PoW/PoS
\end{itemize}
\end{keypointbox}

\subsection{Raft — 非拜占庭容错共识}

\subsubsection{核心思想}

Raft是\textbf{非拜占庭容错}算法，假设所有节点都是诚实的，只处理节点故障（宕机、网络分区），不处理恶意节点。

\subsubsection{三种角色}

\begin{itemize}[leftmargin=1.5em]
\item \textbf{Leader（领导者）}：负责处理所有客户端请求，复制日志到Follower
\item \textbf{Follower（跟随者）}：接收并存储Leader的日志，被动响应
\item \textbf{Candidate（候选人）}：Leader宕机时，Follower变成Candidate，竞选新Leader
\end{itemize}

\subsubsection{Raft vs PBFT}

\begin{center}
\begin{tabular}{|l|p{5cm}|p{5cm}|}
\hline
\rowcolor{LightBG}
\textbf{特性} & \textbf{Raft} & \textbf{PBFT} \\
\hline
容错类型 & 崩溃容错（非恶意） & 拜占庭容错（含恶意节点） \\
\hline
节点假设 & 所有节点诚实 & 最多1/3节点恶意 \\
\hline
通信复杂度 & $O(n)$ & $O(n^2)$ \\
\hline
性能 & 更高（消息少） & 较低（多轮投票） \\
\hline
适用场景 & 私有链、数据库复制 & 联盟链、跨机构协作 \\
\hline
代表项目 & etcd、Consul & Hyperledger Fabric \\
\hline
\end{tabular}
\end{center}

\begin{keypointbox}
重点：Raft vs PBFT的区别：
\begin{itemize}
\item Raft假设节点只是故障，不会作恶（崩溃容错）
\item PBFT假设可能有节点作恶（拜占庭容错）
\item 私有链/内网用Raft，联盟链用PBFT
\end{itemize}
\end{keypointbox}

% =================== 第6章 智能合约 ===================
\section{智能合约}

\subsection{什么是智能合约？}

智能合约（Smart Contract）是\textbf{运行在区块链上的程序代码}，可以自动执行合约条款。

\begin{theorybox}
\textbf{智能合约的核心特点：}
\begin{enumerate}
\item \textbf{自动执行}：满足条件自动触发，无需人工干预
\item \textbf{不可篡改}：一旦部署，代码无法修改
\item \textbf{透明公开}：代码对所有参与者可见
\item \textbf{去中心化}：在区块链网络的所有节点上运行
\end{enumerate}
\end{theorybox}

\subsubsection{生活化类比}

Nick Szabo举了一个经典例子：自动售货机。

\begin{lstlisting}
传统合同：
  1. 你和小明签合同："你给我100元，我给你一个商品"
  2. 你给了100元
  3. 小明可能反悔不给你商品
  4. 你需要去法院起诉

自动售货机（智能合约）：
  1. 你投入100元（满足条件）
  2. 你按下可乐按钮（发出指令）
  3. 机器自动吐出可乐（合同自动执行）
  ——整个过程不需要人工介入，不需要法院介入！
\end{lstlisting}

Nick Szabo认为，法律合同的执行成本很高，而自动化的数字合同可以大幅降低成本。但1994年时没有可信的执行环境，直到以太坊出现。

\subsection{以太坊和EVM}

\subsubsection{以太坊的诞生}

Vitalik Buterin（V神）2013年提出以太坊概念，2015年正式上线。

\textbf{以太坊 = 比特币的区块链 + 图灵完备的虚拟机}

\begin{itemize}[leftmargin=1.5em]
\item 比特币的脚本系统：简单、非图灵完备（不能循环）
\item 以太坊的EVM：图灵完备，可以执行任意复杂逻辑
\end{itemize}

\subsubsection{EVM（以太坊虚拟机）}

EVM是以太坊智能合约的运行环境：
\begin{itemize}[leftmargin=1.5em]
\item 每个以太坊节点都运行EVM
\item 智能合约代码在EVM中执行
\item EVM是\textbf{沙箱环境}，合约不能访问外部网络、文件系统等
\item 合约之间可以互相调用
\end{itemize}

\subsubsection{Gas机制}

执行智能合约需要消耗\textbf{Gas}，防止无限循环和资源滥用。

\begin{lstlisting}
简单转账：21000 Gas（固定）
执行一次算术运算：3 Gas
存储一个32字节数据：20000 Gas
执行一个合约函数：取决于函数复杂度

Gas计算示例：
  用户A调用一个复杂的DeFi合约：
  - 基本操作Gas：50000
  - 存储写入Gas：80000
  - Gas Price：50 Gwei（0.00000005 ETH）
  - 总费用 = 130000 × 50 Gwei = 0.0065 ETH
\end{lstlisting}

\begin{keypointbox}
重点：Gas是防止无限循环和资源滥用的机制。Gas Limit是上限，防止意外超额。EVM（以太坊虚拟机）是智能合约的运行环境，与比特币的脚本系统不同。
\end{keypointbox}

\subsection{Solidity合约示例}

\begin{lstlisting}[language=Solidity]
// SPDX-License-Identifier: MIT
pragma solidity ^0.8.0;

contract SimpleStorage {
    uint256 storedData;
    
    function set(uint256 x) public {
        storedData = x;
    }
    
    function get() public view returns (uint256) {
        return storedData;
    }
}
\end{lstlisting}

调用这个合约：
\begin{lstlisting}
set(42)  ->  交易被广播到网络
           ->  矿工打包到区块
           ->  所有节点执行set函数
           ->  storedData变成42
           ->  状态更新到区块链

get()    ->  本地查询，不消耗Gas
           ->  返回42
\end{lstlisting}

% =================== 第7章 区块链的类型 ===================
\section{区块链的类型}

\subsection{公链（Public Blockchain）}

\begin{itemize}[leftmargin=1.5em]
\item \textbf{特点}：完全开放，任何人可以参与，无需许可
\item \textbf{代表}：比特币、以太坊
\item \textbf{共识}：PoW、PoS
\item \textbf{优势}：去中心化程度高，抗审查，全球可用
\item \textbf{劣势}：性能低，能耗高，监管困难
\end{itemize}

\subsection{联盟链（Consortium Blockchain）}

\begin{itemize}[leftmargin=1.5em]
\item \textbf{特点}：多个机构共同维护，需要许可才能加入
\item \textbf{代表}：Hyperledger Fabric、R3 Corda、FISCO BCOS
\item \textbf{共识}：PBFT、Raft
\item \textbf{优势}：性能高，可监管，适合企业应用
\item \textbf{劣势}：去中心化程度较低
\end{itemize}

\subsection{私链（Private Blockchain）}

\begin{itemize}[leftmargin=1.5em]
\item \textbf{特点}：单一机构控制，完全中心化
\item \textbf{用途}：企业内部审计、数据共享
\item \textbf{优势}：效率最高，完全可控
\item \textbf{劣势}：去中心化程度最低
\end{itemize}

\begin{keypointbox}
重点：联盟链是中国主推方向，在政务、金融、供应链等领域广泛落地。应重点掌握联盟链的PBFT/Raft共识机制和实际应用场景。
\end{keypointbox}

% =================== 第8章 区块链应用领域 ===================
\section{区块链应用领域}

\subsection{数字货币}

\begin{itemize}[leftmargin=1.5em]
\item \textbf{比特币}：数字黄金，价值存储
\item \textbf{稳定币}：USDT、USDC，锚定法币价值
\item \textbf{央行数字货币（CBDC）}：数字人民币（e-CNY）
\end{itemize}

\subsection{供应链金融}

\begin{itemize}[leftmargin=1.5em]
\item 应收账款上链，可追溯、可拆分
\item 降低融资成本，提高资金周转效率
\item 代表：蚂蚁链、微众银行供应链金融平台
\end{itemize}

\subsection{政务与存证}

\begin{itemize}[leftmargin=1.5em]
\item 电子证照上链，防篡改、易验证
\item 司法存证，提高证据可信度
\item 代表：北京互联网法院"天平链"
\end{itemize}

\subsection{数字身份}

\begin{itemize}[leftmargin=1.5em]
\item DID（去中心化身份），用户自主控制
\item 可验证凭证（VC），保护隐私
\item 应用：学历认证、职业资格认证
\end{itemize}

% =================== 第9章 中国区块链发展与国密算法 ===================
\section{中国区块链发展与国密算法}

\subsection{国密算法概述}

中国自主研发的一系列商用密码算法，包括：
\begin{itemize}[leftmargin=1.5em]
\item \textbf{SM2}：非对称加密算法（对应ECDSA/RSA）
\item \textbf{SM3}：Hash算法（对应SHA-256）
\item \textbf{SM4}：对称加密算法（对应AES）
\end{itemize}

\subsection{SM2 — 非对称加密}

基于椭圆曲线密码学，使用sm2p256v1曲线。

\begin{lstlisting}
SM2密钥生成：
1. 生成256位随机数作为私钥
2. 使用椭圆曲线点乘法生成公钥
3. 公钥格式：04 + x坐标(32字节) + y坐标(32字节)

SM2签名过程：
1. 计算消息Hash：e = SM3(message)
2. 生成随机数k
3. 计算椭圆曲线点：(x1, y1) = k * G
4. 计算r = (e + x1) mod n
5. 计算s = ((1 + d)^-1 * (k - r*d)) mod n
6. 签名结果：(r, s)
\end{lstlisting}

\begin{keypointbox}
重点：SM2是中国国密非对称加密算法，与ECDSA对应。如果遇到密码学计算题，关注SM2的签名公式 $(r, s)$ 和验证条件。
\end{keypointbox}

\subsection{SM3 — Hash算法}

输出256位Hash值，与SHA-256安全强度相当。

\begin{lstlisting}
SM3计算示例（简化）：
输入: "Hello SM3"
SM3输出: 66c7f0f462eeedd9d1f2d46b...（64个十六进制字符）

注意：SM3和SHA-256的输出格式相同（64个十六进制字符），
但内部算法不同，输出结果也不同。
\end{lstlisting}

\subsection{SM4 — 对称加密算法}

分组长度为128位，密钥长度128位。

\begin{lstlisting}
SM4加密模式：
- ECB：电子密码本模式（不推荐，相同明文产生相同密文）
- CBC：密码分组链接模式（推荐，需要IV）
- CTR：计数器模式（推荐，可并行加密）

应用场景：
- 数据传输加密
- 存储加密
- 区块链隐私保护（如加密交易数据）
\end{lstlisting}

% =================== 第10章 数据结构基础 ===================
\section{数据结构基础}

\subsection{链表（Linked List）}

\begin{lstlisting}
普通链表：
  节点A -> 节点B -> 节点C -> 节点D
  指针存储下一个节点的内存地址

Hash指针链表（区块链）：
  节点A -> 节点B -> 节点C -> 节点D
  节点A的"指针" = Hash(节点B的内容)
  
  优势：即使有人修改了节点C的数据，
  节点D中存储的Hash(节点C的内容) 立刻失效，
  整个篡改行为被发现。
\end{lstlisting}

\begin{keypointbox}
重点：区块链使用链表结构，每个区块通过Hash指针指向前一个区块。Hash指针 = 数据内容Hash + 位置指针，即使数据移动，Hash值仍然有效，保证链的完整性。
\end{keypointbox}

\subsection{树（Tree）}

\subsubsection{二叉树}

每个节点最多有两个子节点。

\subsubsection{Merkle Tree}

见第3章，区块链核心数据结构。

\subsubsection{MPT（Merkle Patricia Tree）}

以太坊使用的数据结构，结合了：
\begin{itemize}[leftmargin=1.5em]
\item Merkle Tree：验证数据完整性
\item Patricia Trie：高效存储和检索
\end{itemize}

\begin{keypointbox}
重点：以太坊使用MPT而非比特币的简单Merkle Tree，因为MPT支持增删改查操作（而Merkle Tree主要用于验证完整性）。MPT是理解以太坊区块结构的关键。
\end{keypointbox}

% =================== 第11章 链上身份（DID） ===================
\section{链上身份（DID）}

\subsection{什么是DID？}

DID（Decentralized Identifier，去中心化标识符）是\textbf{用户自主控制}的数字身份，不依赖中心化机构（如政府、公司）。

\subsubsection{DID的格式}

\begin{lstlisting}
did:example:123456789abcdefghi
├── did: 方案前缀
├── example: DID方法名（指定操作DID的区块链/网络）
└── 123456789abcdefghi: 唯一标识符（由DID方法生成）
\end{lstlisting}

\subsubsection{DID的生成}

用户通过加密算法生成DID和密钥对，DID本身不包含任何个人身份信息（隐私设计）。

\begin{lstlisting}
DID密钥生成流程：
1. 生成ECDSA secp256k1密钥对（或SM2）
   私钥: 0x7c9e6a5d...
   公钥: 0x04a8e7b6...
2. 计算公钥Hash: SHA256(公钥) -> RIPEMD160 -> Base58
3. 加上方法前缀: did:ethr:0x1234...
4. DID文档存储在以太坊或IPFS上
\end{lstlisting}

\subsubsection{DID文档（DID Document）}

DID文档描述了与DID关联的公钥、服务端点等信息，存储在区块链或去中心化存储上。

\begin{lstlisting}
DID Document示例：
{
  "@context": "https://www.w3.org/ns/did/v1",
  "id": "did:ethr:0x1234abcd...",
  "verificationMethod": [{
    "id": "did:ethr:0x1234#keys-1",
    "type": "EcdsaSecp256k1",
    "controller": "did:ethr:0x1234...",
    "publicKeyHex": "02a8e7b6..."
  }],
  "authentication": ["did:ethr:0x1234...#keys-1"],
  "service": [{
    "type": "LinkedDomains",
    "serviceEndpoint": "https://..."
  }]
}
\end{lstlisting}

DID文档可以更新的操作（更改公钥、服务端点），但DID标识符本身保持不变。

\subsection{可验证凭证（VC）}

VC（Verifiable Credential，可验证凭证）是由权威机构签发的、关于某个主体的数字化声明。

\begin{lstlisting}
VC示例：学历证书
{
  "@context": ["https://www.w3.org/2018/credentials/v1"],
  "type": ["VerifiableCredential", "UniversityDegreeCredential"],
  "issuer": "did:ethr:0xuniversity...",
  "issuanceDate": "2024-06-01",
  "credentialSubject": {
    "id": "did:ethr:0xstudent...",
    "degree": {
      "type": "BachelorDegree",
      "name": "计算机科学",
      "university": "某某大学"
    }
  },
  "proof": {
    "type": "EcdsaSecp256k1Signature2019",
    "created": "2024-06-01",
    "proofValue": "..."
  }
}
\end{lstlisting}

\begin{theorybox}
\textbf{VC的核心特点：}
\begin{enumerate}
\item \textbf{可验证}：任何人都可以验证签发者身份和签名
\item \textbf{防篡改}：任何修改都会导致签名失效
\item \textbf{隐私保护}：可以选择性披露（只证明"我有本科学历"，不透露具体学校）
\item \textbf{自主控制}：用户自己存储和管理凭证，不需要中心数据库
\end{enumerate}
\end{theorybox}

\begin{keypointbox}
重点：DID = 去中心化标识符，由用户自己控制，不依赖中心化机构。三个核心组件：DID标识符、DID文档（存公钥）、可验证凭证（VC）。DID文档存储在区块链或IPFS上。
\end{keypointbox}

% =================== 第12章 实操技能备考要点 ===================
\section{实操技能备考要点}

\subsection{实操技能要点}

本章节主要涵盖实操技能相关知识点：

\begin{enumerate}[leftmargin=1.5em]
    \item 数据结构的基本原理及其在区块链中的应用
    \item 主流密码学算法在区块链身份标识方面的应用
\end{enumerate}

\subsection{Python编程要点}

\subsubsection{Hash计算}

\begin{lstlisting}
import hashlib

# SHA-256计算
def sha256_hash(data):
    if isinstance(data, str):
        data = data.encode('utf-8')
    return hashlib.sha256(data).hexdigest()

# 使用示例
result = sha256_hash("Hello Blockchain")
print(f"SHA-256: {result}")
# 输出: SHA-256: 7f83b1657ff1fc53b92dc18148a1d65dfc2d4b1fa3d677284addd200126d9069
\end{lstlisting}

\subsubsection{Merkle Root计算}

\begin{lstlisting}
import hashlib

def sha256(data):
    if isinstance(data, str):
        data = data.encode('utf-8')
    return hashlib.sha256(data).hexdigest()

def merkle_root(txs):
    """计算Merkle Root"""
    if not txs:
        return sha256("")
    
    # 计算所有交易的Hash
    hashes = [sha256(tx) for tx in txs]
    
    # 逐层计算，直到只剩一个Hash
    while len(hashes) > 1:
        # 如果奇数个，复制最后一个
        if len(hashes) % 2 == 1:
            hashes.append(hashes[-1])
        
        # 两两配对计算父Hash
        parent_hashes = []
        for i in range(0, len(hashes), 2):
            parent = sha256(hashes[i] + hashes[i+1])
            parent_hashes.append(parent)
        
        hashes = parent_hashes
    
    return hashes[0]

# 测试
txs = ["tx1", "tx2", "tx3", "tx4"]
print(f"Merkle Root: {merkle_root(txs)}")
\end{lstlisting}

\subsubsection{简单区块链实现}

\begin{lstlisting}
import hashlib
import time

class Block:
    def __init__(self, index, data, previous_hash):
        self.index = index
        self.timestamp = time.time()
        self.data = data
        self.previous_hash = previous_hash
        self.nonce = 0
        self.hash = self.mine_block()
    
    def calculate_hash(self):
        value = str(self.index) + str(self.timestamp) + \
                str(self.data) + str(self.previous_hash) + \
                str(self.nonce)
        return hashlib.sha256(value.encode()).hexdigest()
    
    def mine_block(self, difficulty=4):
        """PoW挖矿"""
        target = "0" * difficulty
        while True:
            self.hash = self.calculate_hash()
            if self.hash.startswith(target):
                return self.hash
            self.nonce += 1

class Blockchain:
    def __init__(self):
        self.chain = [self.create_genesis_block()]
    
    def create_genesis_block(self):
        return Block(0, "Genesis Block", "0")
    
    def get_latest_block(self):
        return self.chain[-1]
    
    def add_block(self, data):
        previous_hash = self.get_latest_block().hash
        new_block = Block(len(self.chain), data, previous_hash)
        self.chain.append(new_block)
    
    def is_valid(self):
        for i in range(1, len(self.chain)):
            current = self.chain[i]
            previous = self.chain[i-1]
            
            if current.hash != current.calculate_hash():
                return False
            if current.previous_hash != previous.hash:
                return False
        return True

# 使用示例
blockchain = Blockchain()
blockchain.add_block("First Transaction")
blockchain.add_block("Second Transaction")

print(f"区块链有效: {blockchain.is_valid()}")
for block in blockchain.chain:
    print(f"区块{block.index}: {block.hash[:20]}...")
\end{lstlisting}

\subsection{备考建议}

\begin{enumerate}[leftmargin=1.5em]
\item \textbf{理解原理}：不要死记硬背，理解每个算法的原理和应用场景
\item \textbf{动手实践}：至少实现一遍SHA-256、Merkle Tree、简单区块链
\item \textbf{熟悉Python}：掌握基本语法、hashlib库的使用
\item \textbf{编码规范}：代码缩进、命名、注释
\item \textbf{时间分配}：理论45分钟，实操60分钟，合理分配时间
\end{enumerate}

% =================== 附录 ===================
\section{附录A：核心术语速查表}

\begin{center}
\begin{tabular}{|p{3cm}|p{10cm}|}
\hline
\rowcolor{LightBG}
\textbf{术语} & \textbf{解释} \\
\hline
区块链 & 分布式账本技术，数据不可篡改 \\
\hline
Hash & 单向函数，任意输入产生固定长度输出 \\
\hline
SHA-256 & 安全Hash算法，输出256位 \\
\hline
Merkle Tree & 二叉树结构，用于快速验证数据完整性 \\
\hline
UTXO & 未花费交易输出，比特币的记账模型 \\
\hline
PoW & 工作量证明，通过算力竞争记账权 \\
\hline
PoS & 权益证明，通过持币量和时间选择验证者 \\
\hline
PBFT & 实用拜占庭容错，联盟链常用共识 \\
\hline
智能合约 & 运行在区块链上的自动执行程序 \\
\hline
DID & 去中心化身份标识符 \\
\hline
VC & 可验证凭证 \\
\hline
SM2/SM3/SM4 & 中国国密算法 \\
\hline
Gas & 以太坊执行操作的费用单位 \\
\hline
\end{tabular}
\end{center}

\section{附录B：知识体系思维导图}

\begin{theorybox}
\textbf{重点速查（按模块整理）}

\textbf{【必须记忆】}
\begin{itemize}
\item Hash特性：单向、雪崩、抗碰撞
\item SHA-256输出：64个十六进制字符 = 256位
\item ECC vs RSA：256位ECC $\approx$ 3072位RSA
\item 私钥签名、公钥验证
\item 区块头6个字段：Version、PrevHash、Merkle Root、Timestamp、Bits、Nonce
\item UTXO模型：输入引用过去的输出，产生新的输出
\item PoW：算力竞争，找到Nonce使Hash $<$ 目标值
\item PoS：质押代币，随机选择验证者
\item PBFT：三阶段投票，容忍 $<$ 1/3恶意节点
\item 联盟链 vs 公链：PBFT/Raft vs PoW/PoS
\item SM2/SM3/SM4：国密非对称加密、Hash、对称加密
\item DID三要素：标识符、文档、可验证凭证
\end{itemize}

\textbf{【理解即可】}
\begin{itemize}
\item 拜占庭将军问题：叛徒不超过1/3时忠诚将军可达成一致
\item PoS根据持币量和时间选择记账者，不消耗算力
\item Merkle Tree：根Hash代表所有交易，任意一笔修改导致根变化
\item 联盟链用PBFT/Raft，公链用PoW/PoS
\item DID三要素：DID标识符 + DID文档 + 可验证凭证（VC）
\end{itemize}

\textbf{【实操相关】}
\begin{itemize}
\item SHA-256计算：\texttt{hashlib.sha256(data).hexdigest()}
\item 区块链结构：区块头（Hash指针链接）+ 区块体（交易列表）
\item Merkle Root计算：逐层两两Hash，直到只剩一个根Hash
\end{itemize}
\end{theorybox}

\section{附录C：推荐学习资源}

\begin{itemize}[leftmargin=1.5em]
\item \textbf{比特币白皮书}：《Bitcoin: A Peer-to-Peer Electronic Cash System》（中本聪）
\item \textbf{以太坊白皮书}：《A Next-Generation Smart Contract and Decentralized Application Platform》（Vitalik Buterin）
\item \textbf{国密标准}：GM/T 0003-2012（SM2）、GM/T 0004-2012（SM3）、GM/T 0002-2012（SM4）
\item \textbf{在线工具}：\texttt{blockchain.com/explorer}（比特币浏览器）、\texttt{etherscan.io}（以太坊浏览器）
\item \textbf{Python库}：\texttt{hashlib}（标准库）、\texttt{gmssl}（国密算法）
\end{itemize}

\section{附录D：实操环境搭建指南}

\subsection{Python环境准备}

\begin{enumerate}[leftmargin=1.5em]
\item \textbf{安装Python 3.8+}：从官网下载安装 \texttt{python.org}
\item \textbf{安装必要库}：
\begin{lstlisting}
pip install hashlib gmssl pycryptodome
\end{lstlisting}
\item \textbf{验证安装}：
\begin{lstlisting}
import hashlib
print(hashlib.sha256(b"test").hexdigest())
\end{lstlisting}
\end{enumerate}

\subsection{区块链浏览器使用}

\begin{enumerate}[leftmargin=1.5em]
\item \textbf{比特币浏览器}：\texttt{blockchain.com/explorer}
  \begin{itemize}
  \item 输入交易ID查询交易详情
  \item 查看区块高度、时间戳、交易数量
  \end{itemize}
\item \textbf{以太坊浏览器}：\texttt{etherscan.io}
  \begin{itemize}
  \item 查询账户余额和交易历史
  \item 查看智能合约代码
  \end{itemize}
\end{enumerate}

\section{附录E：模拟试题}

\subsection{理论知识题}

\textbf{一、选择题（每题5分）}

\begin{enumerate}[leftmargin=1.5em]
\item SHA-256算法的输出长度是（ ）
  \begin{itemize}
  \item[A] 128位
  \item[B] 256位
  \item[C] 512位
  \item[D] 1024位
  \end{itemize}

\item 比特币使用的共识算法是（ ）
  \begin{itemize}
  \item[A] PoS
  \item[B] PBFT
  \item[C] PoW
  \item[D] Raft
  \end{itemize}

\item PBFT算法能容忍的恶意节点比例最多为（ ）
  \begin{itemize}
  \item[A] 1/2
  \item[B] 1/3
  \item[C] 2/3
  \item[D] 1/4
  \end{itemize}

\item 以下哪个是中国国密非对称加密算法（ ）
  \begin{itemize}
  \item[A] RSA
  \item[B] ECDSA
  \item[C] SM2
  \item[D] AES
  \end{itemize}
\end{enumerate}

\textbf{二、简答题（每题10分）}

\begin{enumerate}[leftmargin=1.5em]
\item 简述区块链的不可篡改性是如何通过密码学实现的。

\item 解释UTXO模型与账户余额模型的区别。

\item 说明Merkle Tree在区块链中的作用。
\end{enumerate}

\subsection{实操技能题}

\textbf{三、编程题（每题20分）}

\begin{enumerate}[leftmargin=1.5em]
\item 编写Python函数计算给定字符串的SHA-256哈希值，并验证雪崩效应。

\textbf{参考答案：}
\begin{lstlisting}
import hashlib

def sha256_hash(data):
    if isinstance(data, str):
        data = data.encode('utf-8')
    return hashlib.sha256(data).hexdigest()

# 验证雪崩效应
h1 = sha256_hash("blockchain")
h2 = sha256_hash("blockchaiN")
print(f"原始: {h1}")
print(f"修改: {h2}")
print(f"不同位数: {sum(c1 != c2 for c1, c2 in zip(h1, h2))}")
\end{lstlisting}

\item 实现简单的Merkle Root计算函数。

\textbf{参考答案：}
\begin{lstlisting}
import hashlib

def sha256(data):
    if isinstance(data, str):
        data = data.encode('utf-8')
    return hashlib.sha256(data).hexdigest()

def merkle_root(txs):
    hashes = [sha256(tx) for tx in txs]
    while len(hashes) > 1:
        if len(hashes) % 2 == 1:
            hashes.append(hashes[-1])
        parent_hashes = []
        for i in range(0, len(hashes), 2):
            parent = sha256(hashes[i] + hashes[i+1])
            parent_hashes.append(parent)
        hashes = parent_hashes
    return hashes[0]

# 测试
txs = ["tx1", "tx2", "tx3", "tx4"]
print(f"Merkle Root: {merkle_root(txs)}")
\end{lstlisting}
\end{enumerate}

\subsection{答案速查}

\begin{center}
\begin{tabular}{|c|c|c|c|}
\hline
\textbf{题号} & \textbf{答案} & \textbf{题号} & \textbf{答案} \\
\hline
1 & B & 3 & B \\
2 & C & 4 & C \\
\hline
\end{tabular}
\end{center}

\end{document}
